\chapter{Cómo reconocer un gato}
% versión completa: chapters/12_recognition.tex

\pencilfig{12}{\pencilart{12}}{Un gato grande y uno pequeño son imágenes muy distintas en los sensores del ojo, pero la respuesta es una sola.}

Un gato en una imagen puede estar a la izquierda o a la derecha, cerca o lejos, al sol o a la sombra. Cada vez cae sobre los sensores del ojo un patrón completamente distinto, y usted lo reconoce en fracciones de segundo. En eso consiste la dificultad del reconocimiento: hay que reconocer lo mismo en miles de aspectos distintos.

\section*{Una cadena de montaje}

El cerebro resuelve el problema con una cadena de montaje de unas diez etapas, y la señal la recorre en quince centésimas de segundo. En cada etapa se hacen dos cosas. Primero, un “y”: la neurona se activa si las piezas necesarias están en el orden necesario, por ejemplo, dos segmentos que se juntan en un ángulo. Después, un “o en todos los lugares”: la neurona siguiente se activa si el ángulo está en cualquier sitio cercano, sin importar dónde exactamente. La primera operación vuelve preciso el reconocimiento; la segunda, tolerante al desplazamiento. Etapa tras etapa, las piezas se juntan en partes, las partes en objetos, y los desplazamientos y las escalas se perdonan.

\pencilfig{112}{%
\foreach \i/\y/\cx/\cy in {1/1.2/0.15/0.3,2/0.3/0.3/0.1,3/-0.6/0.05/0.05}{
  \draw[penlite] (0,\y-0.3) rectangle (0.6,\y+0.3);
  \draw[pcGray, line width=0.8pt] (\cx,\y-0.3+\cy+0.35) -- (\cx,\y-0.3+\cy) -- (\cx+0.3,\y-0.3+\cy);
  \draw[pen=pcBlue, wash=pcBlue] (1.9,\y) circle (0.2);
  \draw[arr] (0.65,\y) -- (1.65,\y);
  \draw[arr=pcBlue] (2.12,\y) -- (3.62,{0.3+0.12*(2-\i)});}
\draw[pen=pcBlue, wash=pcBlue] (3.9,0.3) circle (0.25);
\draw[arr] (4.2,0.3) -- (5.75,0.3); \node[lbl, anchor=south] at (4.97,0.35) {más etapas};
% a small cat
\draw[penlite] (6.3,0.3) circle (0.32);
\draw[penlite] (6.03,0.5) -- (6.07,0.82) -- (6.23,0.6); \draw[penlite] (6.57,0.5) -- (6.53,0.82) -- (6.37,0.6);
\node[lbl, anchor=north] at (0.3,-1.0) {imagen}; \node[lbl, anchor=north, align=center] at (1.9,-1.0) {Y: ángulo\\aquí};
\node[lbl, anchor=north, align=center] at (3.9,-1.0) {O: ángulo\\dondequiera}; \node[lbl, anchor=north] at (6.3,-1.0) {gato};
}{Una etapa de la cadena: el “y” arma el ángulo en un lugar; el “o” perdona dónde está exactamente.}

Justamente este esquema lo copiaron las redes neuronales convolucionales que reconocen imágenes en el teléfono. Y ellas devuelven el favor: las redes entrenadas para reconocer objetos predicen bastante bien cómo responden las neuronas de las etapas superiores de la visión.

\section*{Aprender del video}

Pero ¿quién le enseñó al cerebro? Nadie le mostró a usted un millón de imágenes etiquetadas “gato”. Una respuesta verosímil: el video. El mundo fluye sin cortes, y en una fracción de segundo un objeto casi no cambia; cambian su posición y su iluminación. Si las conexiones se refuerzan entre lo que se ve ahora y lo que se veía un instante antes, la red por sí sola unirá los distintos aspectos de un mismo objeto: lo lento es el objeto; lo rápido, las circunstancias. Hay un experimento en el que los objetos se cambiaban sin que se notara durante los movimientos del ojo, y las neuronas empezaban a tomar objetos distintos por uno solo. Como regla general, esto es todavía una hipótesis.

\section*{Las ilusiones, huellas de los mecanismos}

Las ilusiones no son averías, sino huellas de cómo está hecha la máquina. Si el ojo mira largo rato una cascada y luego unas piedras quietas, las piedras “trepan” hacia arriba: el par de canales “abajo” y “arriba” quedó desequilibrado por el cansancio de uno de ellos. El famoso vestido que unos ven blanco y dorado y otros azul y negro divide a la gente según la iluminación que su cerebro supone en silencio. Las “serpientes que giran” dibujadas parecen moverse porque las zonas brillantes se procesan un poco más rápido que las apagadas, y la diferencia de tiempo se lee como movimiento. Y el cubo que tan pronto sobresale como se hunde son dos grupos de neuronas que compiten entre sí; el cambio de resultado, según un modelo, lo decide sobre todo el ruido.

Es curioso que una red neuronal entrenada para predecir el siguiente fotograma de un video también “ve” girar las serpientes. Así que algunas ilusiones son el precio inevitable de saber predecir.

\section*{El cerebro contra la red neuronal}

El parecido no es completo. Al cerebro le bastan pocos ejemplos, aprende casi sin etiquetas y gasta en todo veinte vatios. A una red neuronal se la puede engañar con una alteración de la imagen que una persona no nota. Eso sí, a las personas también las despistan un poco esas alteraciones si se les muestran por un instante. Dónde pasa la verdadera frontera entre estas dos máquinas todavía se está averiguando.

\fullversion{12}
